\item \model{Kimi K3}

\paragraph*{مقدمه و فلسفه دادگان پیش‌آموزش چندوجهی بومی}
فاز پیش‌آموزش (\en{Pre-training}) مدل \model{Kimi K3} با هدف شکستن مرزهای محاسباتی در مقیاس‌های فراتریلیونی و دستیابی به بافتار فوق‌طولانی ۱ میلیون توکن (\en{1M-token context}) بر پایه معماری تنک ترکیبی از متخصصان (\en{Native Multimodal MoE}) با ۲.۷۸ تریلیون پارامتر کل و ۱۰۴.۲ میلیارد پارامتر فعال طراحی شده است \cite{kimi2026k3}. برخلاف رویکردهای رایج که از یک مرحله انطباق پسینی (\en{Post-hoc Modality Alignment}) جهت اتصال برج بینایی منجمد به مدل زبانی متنی استفاده می‌کنند، در \model{Kimi K3} یک راهبرد یادگیری چندوجهی بومی اتخاذ شده است که در آن توکن‌های متنی و بصری از همان گام نخست آموزش به‌صورت درهم‌تنیده تحت هدف واحد پیش‌بینی توکن بعدی (\en{Next-Token Prediction}) بهینه‌سازی می‌شوند:
\begin{equation}
    \mathcal{L}_{\text{pretrain}}(\theta) = - \sum_{t=1}^{N} \log P_\theta(u_t \mid u_1, u_2, \dots, u_{t-1})
\end{equation}
که در آن $u_t$ نمایانگر یک توکن متنی یا یک توکن بصری فشرده‌شده در فضای برداری مشترک است. ارتقای هم‌زمان خط‌لوله داده، پویایی بهینه‌سازی و بازنویسی متون موجب دستیابی به \textbf{بهره‌وری مقیاس‌بندی ۲.۵ برابری} ($2.5\times$ \en{Scaling Efficiency gain}) در مقایسه با نسل پیشین یعنی \model{Kimi K2} \cite{kimi2025k2} روی داده‌های اعتبارسنجی خارج از توزیع (\en{Held-out OOD}) شده است.

\paragraph*{ترکیب‌بندی و خط لوله پالایش پیکره متنی (\en{Textual Corpus Curation})}
پیکره متنی مدل بر چهار دامنه بنیادین استوار است که با بهره‌گیری از زیرساخت پیشرفته پالایش \cite{kimi2025k2, kimi2026k2_5} پردازش شده‌اند:
\begin{itemize}
    \item \textbf{متن وب (\en{Web Text}):} پوشش‌دهنده تنوع زبانی عمومی و دانش عام حاکم بر بستر اینترنت.
    \item \textbf{کد (\en{Code}):} مخازن سورس‌کد چندزبانه، مستندات فنی و ساختارهای الگوریتمی جهت تقویت توان استدلال صوری.
    \item \textbf{ریاضیات (\en{Mathematics}):} اسناد تحلیلی، متون دارای نمادگذاری علمی لاتک (\en{\LaTeX}) و اثبات‌های ریاضیاتی.
    \item \textbf{دانش تخصصی (\en{Knowledge}):} دانشنامه‌ها، کتب مرجع، اسناد حقوقی، داده‌های مالی و مقالات دانشگاهی بازبینی‌شده.
\end{itemize}

فرایند پاک‌سازی داده‌ها از سه گام دقیق عبور می‌کند:
\begin{enumerate}
    \item \textbf{فیلترهای اکتشافی مبتنی بر قاعده (\en{Rule-based Heuristics}):} حذف محتوای باینری، کدهای خراب، لاگ‌های سیستمی بی‌معنی و متون نامتعارف ماشینی.
    \item \textbf{امتیازدهی کیفی با طبقه‌بند (\en{Classifier-based Quality Scoring}):} تخمین نمره کیفیت اسناد و غربال نمونه‌های با چگالی اطلاعاتی نازل.
    \item \textbf{حذف افزونگی چندمرحله‌ای (\en{Deduplication}):} تلفیق حذف تکرار دقیق با نگاشت‌های قطعی (\en{Exact Dedup}) و حذف تکرار فازی (\en{Fuzzy Dedup}) با تحلیل شباهت سطحی اسناد.
    \item \textbf{نرخ‌های نمونه‌برداری بهینه (\en{Sampling Rates}):} تخصیص بودجه و وزن به هر دامنه بر اساس مطالعات تجربی گسترده بر روی مدل‌های کوچکتری از همان خانواده صورت پذیرفته است.
\end{enumerate}

\paragraph*{دستورالعمل بازنویسی داده‌ها با اعتبارسنجی وفاداری (\en{Fidelity-Verified Rephrasing})}
به منظور برطرف ساختن اطناب کلام و نویزهای ساختاری در متون وب، رویه بازنویسی معرفی‌شده در \model{Kimi K2} \cite{kimi2025k2} با سه راهبرد محوری توسعه یافته است:
\begin{itemize}
    \item \textbf{پرامپت‌نویسی با تنوع دیدگاه و سبک (\en{Style and Perspective-Diverse Prompting}):} متون تخصصی ریاضیاتی و دانشی با زاویه دیدهای متنوع تحلیلی و آموزشی بازآفرینی می‌شوند.
    \item \textbf{تولید خودهمبسته تکه‌ای (\en{Chunk-wise Autoregressive Generation}):} فرایند بازنویسی متن در قالب قطعات منسجم صورت می‌گیرد تا یکدستی ساختار استدلالی تضمین شود.
    \item \textbf{راستی‌آزمایی وفاداری به منبع (\en{Fidelity Verification}):} تمامی داده‌های بازنویسی‌شده توسط ممیزهای خودکار ارزیابی می‌شوند تا با انطباق با سند مرجع، از ورود هرگونه توهم یا داده ساختگی به دادگان پیش‌آموزش جلوگیری شود:
    \begin{equation}
        \text{Fidelity}(D_{\text{rephrased}}, D_{\text{source}}) = \mathbb{I}\left( \text{FactCheck}(D_{\text{rephrased}}, D_{\text{source}}) \ge \tau_{\text{fidelity}} \right)
    \end{equation}
\end{itemize}

\paragraph*{پیکره چندوجهی و آموزش رمزگذار بصری از صفر (\en{MoonViT-V2 from Scratch})}
برخلاف سنت جاری در مدل‌های چندوجهی که رمزگذار تصویر را بر پایه مدل‌های کنتراستیو نظیر \en{SigLIP} مقداردهی اولیه می‌کنند، رمزگذار بینایی \model{Kimi K3} تحت عنوان \model{MoonViT-V2} (با ۴۰۱ میلیون پارامتر و ۲۷ لایه \en{ViT}) به‌طور کامل \textbf{از صفر (\en{From Scratch})} تحت تابع هدف پیش‌بینی توکن بعدی آموزش می‌بیند:
\begin{itemize}
    \item \textbf{پایداری بهینه‌سازی و مهار جهش‌های گرادیان:} آموزش توأم رمزگذارهای کنتراستیو با مدل‌های زبانی عظیم، تلاطم‌های شدیدی در نُرم گرادیان برج بینایی پدید می‌آورد. همان‌طور که در شواهد گزارش فنی مشخص است، مدل پیشین با مقداردهی \en{SigLIP} دچار جهش‌های مقطعی مکرر در نُرم گرادیان می‌شد، درحالی‌که \model{MoonViT-V2} نُرم‌های گرادیان بسیار پایین‌تر و کاملاً پایداری را تجربه می‌کند.
    \item \textbf{داده‌های برنامه‌نویسی بصری (\en{Programmatic Multimodal Data}):} علاوه بر داده‌های کلاسیک تصویر-زیرنویس، اسناد متنی-تصویری درهم‌تنیده و \en{OCR}، حجم عظیمی از داده‌های برنامه‌نویسی شامل کدهای منبع به همراه خروجی‌های رندرشده دوبعدی و سه‌بعدی (\en{SVG}، صفحات وب، بازی‌های رایانه‌ای، دارایی‌های سه‌بعدی و نقشه‌های مهندسی \en{CAD}) به پیکره آموزشی تزریق گردید.
    \item \textbf{نظارت مختصات فضایی:} تمام نمونه‌های نیازمند درک مکان به دو صورت مختصات مطلق پیکسلی و نرمال‌شده در بازه $[0, 1]^2$ برچسب‌گذاری شدند.
    \item \textbf{فشرده‌سازی پیکسلی توکن‌ها (\en{Pixel-Shuffle Downsampling}):} تصاویر و فریم‌های ویدئویی تا وضوح بالای $3584 \times 3584$ با کاهش نمونه‌برداری $2 \times 2$ از طریق عملگر \en{Pixel-Shuffle} متراکم شده و تعداد توکن‌های بینایی را با ضریب ۴ کاهش می‌دهند تا گنجایش بافت ۱ میلیون توکنی حفظ شود:
    \begin{equation}
        T_{\text{visual}} = \frac{1}{4} \cdot \left( \frac{H_{\text{img}}}{P} \times \frac{W_{\text{img}}}{P} \right), \quad P = 14
    \end{equation}
\end{itemize}

\paragraph*{مهندسی داده‌های بافتار فوق‌بلند و مسائل مصنوعی}
به منظور تثبیت ظرفیت بازیابی و استدلال در سقف بافت ۱ میلیون توکن:
\begin{itemize}
    \item \textbf{پالایش ویدئوها با هش ادراکی (\en{Perceptual Hashing}):} فریم‌های ویدئویی تکراری و ساکن با استفاده از \en{pHash} غربال می‌شوند تا از افزونگی بیهوده توکن‌های بصری جلوگیری گردد.
    \item \textbf{وزن‌دهی افزایشی در فاز سردسازی (\en{Upsampling during Cooldown}):} به سبب کمیابی طبیعی اسناد طولانی منسجم در وب، فراوانی اسناد طولانی و ویدئوها در مراحل پایانی پیش‌آموزش و به‌ویژه در فاز سردسازی (\en{Cooldown Phase}) به‌طور چشمگیری تقویت گردید.
    \item \textbf{تولید وظایف ترکیبی با وابستگی سراسری (\en{Synthetic Interleaved Tasks}):} برای جلوگیری از فروپاشی سازوکارهای توجه به الگوهای محلی، اسناد و زیروظایف گوناگون به‌گونه‌ای چیده و ادغام شدند که حل مسئله نهایی منوط به تلفیق اطلاعاتی باشد که در سراسر بازه ۱ میلیون توکن پراکنده‌اند:
    \begin{equation}
        y_T = f(x_{t_1}, x_{t_2}, \dots, x_{t_k}), \quad 0 \ll t_1 \ll t_2 \ll \dots \ll t_k \le 10^6
    \end{equation}
\end{itemize}

\paragraph*{برنامه درسی پیشرونده طول بافتار (\en{Progressive Context Curriculum})}
مدیریت محاسبات پیش‌آموزش بر روی طول توالی‌های بالا از طریق یک برنامه درسی چهارمرحله‌ای انجام گرفته است:
\begin{itemize}
    \item \textbf{مرحله اول (پیش‌آموزش پایه):} آموزش در طول بافت $8\text{K}$ توکن آغاز می‌شود.
    \item \textbf{مرحله دوم (گسترش میانی):} افزایش پنجره بافتار به $64\text{K}$ توکن در ادامه فرایند پیش‌آموزش.
    \item \textbf{مرحله سوم و چهارم (فاز سردسازی):} جهش متوالی طول بافت از $256\text{K}$ به $1\text{M}$ توکن در طول دوره \en{Cooldown}.
\end{itemize}
با اتکا به مکانیزم زوال ذاتی لایه‌های توجه خطی \en{KDA} \cite{kimi2025linear} و حذف کامل کدهای موضعی صریح (\en{NoPE}) در لایه‌های \en{Gated MLA}، مدل بدون نیاز به دستکاری پایه‌های فرکانسی \en{RoPE} یا روش‌های برون‌یابی نظیر \en{YaRN} \cite{peng2023yarn}، مستقیماً به بافتار ۱ میلیون توکنی تعمیم می‌یابد.

\paragraph*{قوانین مقیاس‌پذیری داده‌ها و برنامه واپاشی نرخ یادگیری}
روابط تجربی حاکم بر مصرف داده و توان پردازشی بر پایه قوانین توان (\en{Power-law Scaling}) تدوین شده‌اند:
\begin{equation}
    \mathcal{L}_{\text{val}}(C) = \mathcal{L}_{\infty} + \left(\frac{C_0}{C}\right)^{\alpha}
\end{equation}
برازش منحنی‌های قانون مقیاس‌پذیری حاکی از آن است که \model{Kimi K3} برای دستیابی به تراز زیان اعتبارسنجی مشابه با \model{Kimi K2}، به تقریب تنها به ۴۰٪ از حجم پردازش ($C_{\text{K3}} \approx 0.4 \cdot C_{\text{K2}}$) نیاز دارد. 

علاوه بر این، در مقایسه زمان‌بند واپاشی کسینوسی (\en{Cosine Decay}) با زمان‌بند گرمایش-پایداری-زوال (\en{WSD}) \cite{hu2024minicpm}، نشان داده شد که با اجرای مستقل جستجوی قوانین مقیاس‌پذیری جهت تعیین نرخ یادگیری اوج و اندازه دسته بهینه برای هر روش، رویکرد \textbf{واپاشی کسینوسی با ۱٪ گرمایش خطی} به طور مداوم زیان نهایی پایین‌تری نسبت به \en{WSD} ثبت کرده و به عنوان زمان‌بند رسمی پیش‌آموزش تثبیت شد. پارامترهای ماتریسی مدل نیز با بهینه‌ساز تفکیک‌شده \en{Per-Head Muon} \cite{jordan2024muon} و زوال وزن $0.1$ آموزش داده شدند.

\paragraph*{مشخصات آماری و معماری داده‌محور (\en{Architectural Specifications})}
جدول \ref{tab:kimi_k3_arch_specs} مقایسه مشخصات و پارامترهای پیش‌آموزش مدل \model{Kimi K3} را در برابر \model{Kimi K2} نشان می‌دهد. این مشخصات ظرفیت داده‌پذیری شبکه را در سطح ۲.۷۸ تریلیون پارامتر با ۸۹۶ متخصص مسیردهی‌شده تبیین می‌کند.

\begin{table}[htbp]
\centering
\caption{مقایسه معماری و مشخصات پیش‌آموزش مدل‌های \model{Kimi K2} و \model{Kimi K3}}
\label{tab:kimi_k3_arch_specs}
\resizebox{\textwidth}{!}{%
\begin{tabular}{lccc}
\toprule
\textbf{ویژگی / مؤلفه معماری} & \textbf{\model{Kimi K2}} & \textbf{\model{Kimi K3}} & \textbf{تغییرات ($\Delta$)} \\
\midrule
معماری کلی (\en{Architecture}) & \en{MoE} & \en{MoE} & — \\
تعداد کل لایه‌ها (\en{\#Layers}) & 61 & 93 & $\uparrow 52\%$ \\
تعداد کل پارامترها (\en{Total Parameters}) & 1.04T & \textbf{2.78T} & $\uparrow 167\%$ \\
پارامترهای فعال به‌ازای هر توکن (\en{Activated Parameters}) & 32.6B & \textbf{104.2B} & $\uparrow 220\%$ \\
ابعاد پنهان لایه‌ها (\en{Hidden Dimension}) & 7,168 & 7,168 & $=$ \\
ابعاد لایه فشرده نهفته (\en{Latent MoE Dimension}) & — & 3,584 ($0.5\times$) & — \\
ابعاد پنهان به‌ازای هر متخصص (\en{MoE Hidden Dim per Expert}) & 2,048 & 3,072 & $\uparrow 50\%$ \\
تعداد کل متخصصان مسیردهی‌شده (\en{Routed Experts}) & 384 & \textbf{896} & $\uparrow 133\%$ \\
تعداد متخصصان فعال در هر توکن (\en{Active Experts per Token}) & 8 & \textbf{16} & $\uparrow 100\%$ \\
تعداد متخصصان اشتراکی (\en{Shared Experts}) & 1 & 2 & $\uparrow 100\%$ \\
تعداد سرهای توجه (\en{Attention Heads}) & 64 & 96 & $\uparrow 50\%$ \\
تعداد لایه‌های متراکم (\en{Dense Layers}) & 1 & 1 & $=$ \\
اندازه دایره واژگان (\en{Vocabulary Size}) & 160K & 160K & $=$ \\
طول بافتار پیش‌آموزش (\en{Training Context Length}) & 128K & \textbf{1M} & $\mathbf{\uparrow 8\times}$ \\
مکانیزم توجه (\en{Attention Mechanism}) & \en{MLA} & \textbf{\en{Hybrid KDA--MLA}} & — \\
ترکیب‌بندی لایه‌های توجه & 61 \en{MLA} & 69 \en{KDA} + 24 \en{MLA} & — \\
تابع فعال‌ساز (\en{Activation Function}) & \en{SwiGLU} \cite{shazeer2020glu} & \textbf{\en{SiTU-GLU}} & — \\
تعداد لایه‌های پیش‌بینی چندتوکنی (\en{MTP Layers}) & 1 layer & 1 layer & $=$ \\
کل پارامترهای برج بینایی (\en{ViT Total Parameters}) & — & \textbf{401M} & — \\
تعداد لایه‌های برج بینایی (\en{\#ViT Layers}) & — & 27 layers & — \\
اندازه پچ ورودی تصویر (\en{Patch Size}) & — & 14 & — \\
تعداد سرهای توجه برج بینایی (\en{ViT Attention Heads}) & — & 12 & — \\
\bottomrule
\end{tabular}%
}
\end{table}